\clearpage
\section{Memory Model}

The memory model defines instruction fetches, loads, stores, atomic read-modify-write operations, cache maintenance,
and translation-cache maintenance after EA calculation and address translation.

\subsection{Baseline Ordering}
Normal memory is coherent. Every hardware thread observes writes to a given byte in one common coherence order, and
writes by one hardware thread to that byte enter the order in program order. A naturally aligned tear-free scalar
store enters that order as one write event.

Normal-memory writes are multi-copy atomic. When a write event becomes observable by a hardware thread other than
the writer, it becomes observable to all such threads at the same architectural point. An implementation may still
forward a hardware thread's own pending store to its later loads before that store is globally observable.

The baseline ordering is weak. Except for same-byte coherence or ordering imposed by an atomic memory-order selector
or fence, loads and stores to different locations may be issued, completed, and observed in an order different from
program order. Coherence and multi-copy atomicity do not by themselves order accesses to different locations.
Platform-defined device memory may impose additional rules outside this normal-memory model.

\subsection{Access Atomicity and Alignment}
Unaligned ordinary memory accesses are architecturally permitted and may be implemented as multiple aligned memory
beats. A successfully completed unaligned access has no tear-free guarantee; each visible store beat is separately
coherent and multi-copy atomic.

An unaligned store is nevertheless fault-atomic for synchronous faults. If any byte of the complete store range
would raise a synchronous fault, no byte of the store becomes architecturally visible. An implementation that uses
multiple beats must validate, buffer, or roll back those beats to preserve this rule. This fault atomicity does not
make a successfully completed unaligned store tear-free.

Atomic read-modify-write instructions require naturally aligned memory operands. A misaligned atomic operand raises
\texttt{PAGE\_FAULT} with the \texttt{ALIGNMENT\_ATOMIC} subtype before any memory update or architectural result
commits; it must not complete as a non-atomic update.

Stores to memory that is later executed as instructions are ordinary data stores until an explicit synchronization
sequence makes the modified bytes visible to instruction fetch.

For every scalar size below, a naturally aligned normal-memory load or store is tear-free. Unaligned accesses retain
the successful-access and synchronous-fault rules above.

\Needspace{2.3in}
\manualtablecaption{Scalar Access Sizes and Alignment}
\begin{manualtabular}{@{}>{\raggedright\arraybackslash}p{0.75in}>{\raggedright\arraybackslash}p{0.70in}>{\raggedright\arraybackslash}p{3.95in}@{}}
\toprule
\rowcolor{ManualHeaderFill}
\textbf{Suffix} & \textbf{Bytes} & \textbf{Natural Alignment}\\
\midrule
B & 1 & any byte address\\
W & 2 & address divisible by 2\\
L & 4 & address divisible by 4\\
Q & 8 & address divisible by 8\\
\bottomrule
\end{manualtabular}

\subsection{Atomic Memory-Order Selectors}
\manualtablecaption{Atomic Memory-Order Selectors}
\begin{manualtabular}{@{}>{\raggedright\arraybackslash}p{0.95in}>{\raggedright\arraybackslash}p{0.45in}>{\raggedright\arraybackslash}p{4.00in}@{}}
\toprule
\rowcolor{ManualHeaderFill}
\textbf{Selector} & \textbf{Code} & \textbf{Architectural Ordering Effect}\\
\midrule
\texttt{RELAXED} & \texttt{0} & Atomicity only. No ordering is imposed on unrelated memory operations.\\
\texttt{ACQUIRE} & \texttt{1} & Later memory operations by the same hardware thread are ordered after the atomic operation.\\
\texttt{RELEASE} & \texttt{2} & Earlier memory operations by the same hardware thread are ordered before the atomic operation.\\
\texttt{ACQREL} & \texttt{3} & Applies both acquire and release ordering.\\
\texttt{SEQCST} & \texttt{4} & Applies acquire-release ordering and participates in the global sequentially consistent order for atomic operations.\\
\texttt{RESERVED5} & \texttt{5} & Reserved selector encoding; using it is invalid for the atomic form.\\
\texttt{RESERVED6} & \texttt{6} & Reserved selector encoding; using it is invalid for the atomic form.\\
\texttt{RESERVED7} & \texttt{7} & Reserved selector encoding; using it is invalid for the atomic form.\\
\bottomrule
\end{manualtabular}

\subsection{Fence Instructions}
\hyperref[instr:rfence]{\texttt{RFENCE}} orders prior reads before later reads,
\hyperref[instr:wfence]{\texttt{WFENCE}} orders prior writes before later writes, and
\hyperref[instr:afence]{\texttt{AFENCE}} orders prior memory operations before later memory operations issued by
the same hardware thread. The linked instruction entries define the normative ordering behavior.
